\documentclass[11pt,a4paper]{article}

% ---------------------------------------------------------
% PACKAGES
% ---------------------------------------------------------
\usepackage[utf8]{inputenc}
\usepackage[T1]{fontenc}
\usepackage{lmodern}
\usepackage{amsmath, amssymb, amsfonts}
\usepackage{bm}
\usepackage{graphicx}
\usepackage{hyperref}
\usepackage{authblk}
\usepackage{geometry}
\usepackage{cite}
\usepackage{physics}
\usepackage{tensor}
\usepackage{enumitem}

\geometry{margin=1in}

% ---------------------------------------------------------
% TITLE
% ---------------------------------------------------------
\title{\textbf{Baryonic Physics and Feedback in the Relaxation‑Driven Cyclic Cosmology  
Cooling, Star Formation, Outflows, and the RDCC Baryon Connection}}

\author[1]{Michael Lehmann}
\affil[1]{\small Independent Researcher, Relaxation-Driven Cyclic Cosmology (RDCC) Framework\\
\texttt{mi.lehmann@gmx.de}}

\date{\small \today}

% ---------------------------------------------------------
% DOCUMENT START
% ---------------------------------------------------------
\begin{document}

\maketitle

\begin{abstract}
Baryonic physics plays a decisive role in shaping the observable Universe. Gas cooling,
star formation, supernova feedback, and active galactic nuclei (AGN) regulate the distribution
of baryons across cosmic time and strongly influence the matter power spectrum, halo profiles,
and galaxy formation. In the standard $\Lambda$CDM model, baryonic feedback is typically
treated as an astrophysical correction layered on top of scale-independent gravitational collapse.
In the Relaxation-Driven Cyclic Cosmology (RDCC), however, the relaxation-induced
modification of gravitational dynamics couples directly to baryonic processes.

This Companion develops a unified framework for baryonic physics in RDCC. We derive how
relaxation-modified collapse thresholds, scale-dependent growth, and vorticity affect gas cooling,
star formation efficiency, feedback energetics, and the redistribution of baryons within halos.
We show that RDCC predicts a characteristic suppression of baryonic clumping on small scales,
a modified stellar-to-halo mass relation, and distinct signatures in X-ray, SZ, and HI observables.

Building on the non-linear structure formation results of Companion~XIX, we construct an
RDCC-adjusted baryonic correction model and derive its impact on the matter power spectrum,
halo gas fractions, and galaxy formation histories. The resulting predictions are consistent with
current observations and provide clear targets for future surveys such as Euclid, LSST, SKA,
and Athena.

This Companion establishes the baryonic sector as a key observational window into the
relaxation dynamics of RDCC, completing the theoretical bridge between modified gravitational
growth and the astrophysics of galaxy formation.
\end{abstract}

% ---------------------------------------------------------
% SECTION 1 — THE RDCC FRAMEWORK AND ITS IMPLICATIONS FOR BARYONIC PHYSICS
% ---------------------------------------------------------
\section{The RDCC Framework and Its Implications for Baryonic Physics}

The Relaxation-Driven Cyclic Cosmology (RDCC) modifies the growth of cosmic structures
through a scale-dependent relaxation rate $\alpha(a)$, which introduces a characteristic
comoving relaxation scale $k_{\mathrm{rel}}(a)$ and an associated mass scale $M_{\mathrm{rel}}(a)$.
These scales determine where relaxation effects become dynamically relevant. Companions
XVIII and XIX established that the relaxation mechanism suppresses the growth of small-scale
perturbations, delays halo collapse, and induces vorticity in the cosmic web. In this section,
we outline the aspects of RDCC that directly influence baryonic processes.

Relaxation modifies baryonic physics through:
\begin{itemize}
\item reduced gravitational potentials in low-mass halos,
\item delayed collapse and lower virial temperatures,
\item reduced inflow velocities and enhanced vorticity,
\item increased sensitivity to feedback-driven gas loss,
\item smoother gas distributions in the CGM and IGM.
\end{itemize}

These effects form a coherent physical picture in which relaxation dynamics and baryonic
processes are tightly coupled.

% ---------------------------------------------------------
\subsection{Relaxation-Modified Gravitational Collapse}

In $\Lambda$CDM, the collapse of halos is governed by a scale-independent linear growth
function. In RDCC, however, modes with $k > k_{\mathrm{rel}}$ experience suppressed growth,
leading to delayed formation of low-mass halos. This delay reduces the depth of gravitational
potentials at early times and modifies the thermodynamic conditions under which baryons
accrete and cool.

The virial temperature of a halo in RDCC can be approximated as
\begin{equation}
T_{\mathrm{vir,RDCC}}
=
T_{\mathrm{vir},\Lambda\mathrm{CDM}}
\left[
1
-
\chi_{T}
\left(
\frac{M_{\mathrm{rel}}}{M}
\right)^{\beta_{T}}
\right],
\end{equation}
with typical values $\chi_{T} \sim 0.05$ and $\beta_{T} \sim 2/3$.

Lower virial temperatures imply:
\begin{itemize}
\item reduced cooling efficiencies,
\item delayed formation of cold gas disks,
\item suppressed early star formation,
\item a shift of the cosmic star formation peak to lower redshift.
\end{itemize}

% ---------------------------------------------------------
\subsection{Scale-Dependent Growth and Vorticity}

A distinctive feature of RDCC is the emergence of relaxation-induced vorticity. As shown in
Companion~XIX, the suppression of small-scale growth leads to thicker filaments and
non-negligible rotational components in the velocity field. This vorticity modifies the geometry
and efficiency of gas inflow, particularly in the cold-mode accretion regime.

The inflow velocity becomes
\begin{equation}
v_{\mathrm{inflow,RDCC}}
=
v_{\mathrm{inflow},\Lambda\mathrm{CDM}}
\left[
1
-
\chi_{v}
\left(
\frac{k}{k_{\mathrm{rel}}}
\right)^{\alpha}
\right],
\end{equation}
with $\chi_{v} \sim 0.1$.

Consequences include:
\begin{itemize}
\item reduced gas supply at early times,
\item thicker and more turbulent filaments,
\item lower central gas densities,
\item modified angular momentum acquisition.
\end{itemize}

% ---------------------------------------------------------
\subsection{Relaxation Mass Scale and Baryonic Sensitivity}

The relaxation mass scale $M_{\mathrm{rel}}$ determines which halos experience significant
baryonic modifications. Halos with $M < M_{\mathrm{rel}}$:
\begin{itemize}
\item form later,
\item have shallower potentials,
\item are more susceptible to feedback-driven gas loss,
\item exhibit reduced cooling and star formation.
\end{itemize}

Halos with $M > M_{\mathrm{rel}}$ behave similarly to their $\Lambda$CDM counterparts, with
only mild relaxation effects.

This mass dependence naturally explains:
\begin{itemize}
\item the baryon deficit in dwarf galaxies,
\item reduced stellar-to-halo mass ratios in low-mass systems,
\item smoother gas profiles in galaxy groups,
\item the robustness of cluster-scale baryon fractions.
\end{itemize}

% ---------------------------------------------------------
\subsection{Implications for Baryonic Processes}

The relaxation-modified gravitational environment affects baryonic physics in four key ways:

\paragraph{1. Gas cooling.}
Reduced virial temperatures and delayed collapse increase cooling times and suppress the
formation of cold gas.

\paragraph{2. Star formation.}
Lower gas densities and longer dynamical times reduce star formation efficiency, especially
in halos with $M < M_{\mathrm{rel}}$.

\paragraph{3. Feedback.}
Shallower potentials enhance the ability of SN and AGN feedback to expel gas, increasing
mass-loading factors and redistributing baryons to larger radii.

\paragraph{4. Baryon redistribution.}
Relaxation-induced vorticity and reduced inflow reshape the circumgalactic and intergalactic
media, producing smoother gas profiles and more extended CGM structures.

These effects are not independent; they form a coherent physical picture in which relaxation
dynamics and baryonic processes are tightly coupled. The following sections develop each of
these aspects in detail.

% ---------------------------------------------------------
% SECTION 2 — GAS COOLING AND STAR FORMATION IN THE RDCC FRAMEWORK
% ---------------------------------------------------------
\section{Gas Cooling and Star Formation in the RDCC Framework}

Gas cooling and star formation are among the most sensitive baryonic processes affected by
relaxation-modified gravitational dynamics. In $\Lambda$CDM, these processes are primarily
governed by the depth of the gravitational potential, the thermodynamic state of the gas, and
the efficiency of radiative cooling. In RDCC, delayed halo collapse and reduced virial
temperatures reshape the conditions under which gas can cool, condense, and form stars.
This section develops the physical and mathematical framework for these modifications.

Relaxation affects cooling and star formation through:
\begin{itemize}
\item reduced virial temperatures in halos with $M < M_{\mathrm{rel}}$,
\item reduced inflow velocities due to relaxation-induced vorticity,
\item longer cooling times and delayed disk formation,
\item reduced star formation efficiency and suppressed early SFR,
\item lower cold-gas fractions and smoother gas distributions.
\end{itemize}

% ---------------------------------------------------------
\subsection{Cooling Times and Virial Temperatures}

The radiative cooling time of gas with temperature $T$ and electron density $n_{e}$ is
\begin{equation}
t_{\mathrm{cool}}
=
\frac{3 k_{\mathrm{B}} T}{2 n_{e} \Lambda(T,Z)},
\end{equation}
where $\Lambda(T,Z)$ is the cooling function.

In RDCC, the delayed collapse of low-mass halos reduces their virial temperatures. Using the
scaling from Section~1,
\begin{equation}
T_{\mathrm{vir,RDCC}}
=
T_{\mathrm{vir},\Lambda\mathrm{CDM}}
\left[
1
-
\chi_{T}
\left(
\frac{M_{\mathrm{rel}}}{M}
\right)^{\beta_{T}}
\right],
\end{equation}
the cooling time becomes
\begin{equation}
t_{\mathrm{cool,RDCC}}
=
t_{\mathrm{cool},\Lambda\mathrm{CDM}}
\left[
1
+
\chi_{\mathrm{cool}}
\left(
\frac{M_{\mathrm{rel}}}{M}
\right)^{\beta_{\mathrm{cool}}}
\right],
\end{equation}
with $\chi_{\mathrm{cool}} \sim 0.1$.

Physical consequences:
\begin{itemize}
\item reduced cooling rates in halos with $M \lesssim M_{\mathrm{rel}}$,
\item delayed formation of cold gas disks,
\item suppression of early star formation,
\item a shift of the cosmic star formation peak to lower redshift.
\end{itemize}

% ---------------------------------------------------------
\subsection{Filamentary Inflow and Gas Supply}

Gas accretion onto halos occurs through two primary channels:
\begin{enumerate}
\item cold-mode accretion along filaments,
\item hot-mode accretion through virial shocks.
\end{enumerate}

Relaxation-induced vorticity modifies the geometry and efficiency of cold-mode inflow. The
reduced inflow velocity is
\begin{equation}
v_{\mathrm{inflow,RDCC}}
=
v_{\mathrm{inflow},\Lambda\mathrm{CDM}}
\left[
1
-
\chi_{v}
\left(
\frac{k}{k_{\mathrm{rel}}}
\right)^{\alpha}
\right],
\end{equation}
with $\chi_{v} \sim 0.1$.

Consequences:
\begin{itemize}
\item thicker filaments,
\item lower gas supply at early times,
\item reduced central gas densities,
\item modified angular momentum acquisition.
\end{itemize}

% ---------------------------------------------------------
\subsection{Star Formation Efficiency}

The star formation rate (SFR) is commonly modeled as
\begin{equation}
\dot{M}_{\star}
=
\epsilon_{\star}
\frac{M_{\mathrm{gas}}}{t_{\mathrm{dyn}}},
\end{equation}
where $\epsilon_{\star}$ is the star formation efficiency and $t_{\mathrm{dyn}}$ is the dynamical time.

In RDCC, both $M_{\mathrm{gas}}$ and $t_{\mathrm{dyn}}$ are modified:
\begin{itemize}
\item reduced inflow lowers $M_{\mathrm{gas}}$,
\item delayed collapse increases $t_{\mathrm{dyn}}$,
\item shallower potentials reduce gas compression.
\end{itemize}

The resulting star formation efficiency becomes
\begin{equation}
\epsilon_{\star,\mathrm{RDCC}}
=
\epsilon_{\star,\Lambda\mathrm{CDM}}
\left[
1
-
\chi_{\star}
\left(
\frac{M_{\mathrm{rel}}}{M}
\right)^{\beta_{\star}}
\right],
\end{equation}
with $\chi_{\star} \sim 0.1$.

Consequences:
\begin{itemize}
\item star formation in low-mass halos is strongly suppressed,
\item the cosmic SFR peak shifts by $\Delta z \sim 0.3$ to lower redshift,
\item stellar-to-halo mass ratios are reduced in dwarf galaxies.
\end{itemize}

% ---------------------------------------------------------
\subsection{Cold Gas Fractions}

The cold gas fraction is defined as
\begin{equation}
f_{\mathrm{cold}}
=
\frac{M_{\mathrm{cold}}}{M_{\mathrm{gas}}}.
\end{equation}

In RDCC, the reduced inflow and longer cooling times lead to
\begin{equation}
f_{\mathrm{cold,RDCC}}
=
f_{\mathrm{cold},\Lambda\mathrm{CDM}}
\left[
1
-
\chi_{\mathrm{cold}}
\left(
\frac{M_{\mathrm{rel}}}{M}
\right)^{\beta_{\mathrm{cold}}}
\right],
\end{equation}
with $\chi_{\mathrm{cold}} \sim 0.05$.

Predictions:
\begin{itemize}
\item lower HI masses in dwarf galaxies,
\item reduced molecular gas fractions,
\item delayed disk formation,
\item smoother gas distributions in low-mass halos.
\end{itemize}

% ---------------------------------------------------------
\subsection{Cosmic Star Formation History}

The cosmic star formation rate density is
\begin{equation}
\dot{\rho}_{\star}(z)
=
\int dM\,
\dot{M}_{\star}(M,z)\,
n(M,z),
\end{equation}
where $n(M,z)$ is the halo mass function.

Using the RDCC-modified halo mass function and SFR, we find:
\begin{itemize}
\item a suppressed high-redshift SFR,
\item a delayed SFR peak,
\item a mildly enhanced low-redshift tail.
\end{itemize}

These trends are consistent with JWST observations of early galaxy populations.

% ---------------------------------------------------------
% SECTION 3 — SUPERNOVA AND AGN FEEDBACK IN THE RELAXATION-MODIFIED REGIME
% ---------------------------------------------------------
\section{Supernova and AGN Feedback in the Relaxation-Modified Regime}

Feedback from supernovae (SN) and active galactic nuclei (AGN) plays a decisive role in
regulating star formation and redistributing baryons within halos. In $\Lambda$CDM, the
efficiency of feedback is primarily determined by the depth of the gravitational potential and
the density structure of the interstellar and circumgalactic media. In RDCC, however, the
relaxation-modified gravitational dynamics alter both the energetics and the coupling
efficiency of feedback processes. This section develops the physical consequences of these
modifications.

Relaxation affects feedback through:
\begin{itemize}
\item reduced escape velocities in halos with $M < M_{\mathrm{rel}}$,
\item reduced central gas densities due to suppressed inflow,
\item enhanced propagation of SN- and AGN-driven outflows,
\item increased coupling efficiency of AGN energy,
\item more extended redistribution of baryons into the CGM and IGM.
\end{itemize}

% ---------------------------------------------------------
\subsection{Supernova Feedback and Escape Velocities}

The total energy injected by supernovae associated with a stellar mass $\Delta M_{\star}$ is
\begin{equation}
E_{\mathrm{SN}}
=
\epsilon_{\mathrm{SN}}
\Delta M_{\star}
\times 10^{51}\,\mathrm{erg},
\end{equation}
where $\epsilon_{\mathrm{SN}}$ is the number of supernovae per solar mass of stars formed.

The ability of SN feedback to expel gas depends on the escape velocity,
\begin{equation}
v_{\mathrm{esc}}
=
\sqrt{2|\Phi_{\mathrm{halo}}|},
\end{equation}
where $\Phi_{\mathrm{halo}}$ is the gravitational potential.

In RDCC, the delayed collapse of halos reduces the depth of the potential well:
\begin{equation}
\Phi_{\mathrm{RDCC}}
=
\Phi_{\Lambda\mathrm{CDM}}
\left[
1
-
\chi_{\Phi}
\left(
\frac{M_{\mathrm{rel}}}{M}
\right)^{\beta_{\Phi}}
\right],
\end{equation}
with $\chi_{\Phi} \sim 0.1$.

Consequently, the escape velocity becomes
\begin{equation}
v_{\mathrm{esc,RDCC}}
=
v_{\mathrm{esc},\Lambda\mathrm{CDM}}
\left[
1
-
\chi_{v,\mathrm{esc}}
\left(
\frac{M_{\mathrm{rel}}}{M}
\right)^{\beta_{v}}
\right],
\end{equation}
with $\chi_{v,\mathrm{esc}} \sim 0.1$.

Implications:
\begin{itemize}
\item SN-driven outflows are more efficient in low-mass halos,
\item gas removal is enhanced, reducing central gas densities,
\item metal enrichment of the circumgalactic medium (CGM) increases.
\end{itemize}

% ---------------------------------------------------------
\subsection{Mass-Loading Factors}

The mass-loading factor, defined as
\begin{equation}
\eta
=
\frac{\dot{M}_{\mathrm{out}}}{\dot{M}_{\star}},
\end{equation}
quantifies the efficiency of outflows.

In RDCC, the enhanced efficiency of SN-driven winds leads to
\begin{equation}
\eta_{\mathrm{RDCC}}
=
\eta_{\Lambda\mathrm{CDM}}
\left[
1
+
\chi_{\eta}
\left(
\frac{M_{\mathrm{rel}}}{M}
\right)^{\beta_{\eta}}
\right],
\end{equation}
with $\chi_{\eta} \sim 0.1$.

Predictions:
\begin{itemize}
\item stronger outflows in dwarf galaxies,
\item reduced baryon fractions in low-mass halos,
\item enhanced metal transport into the IGM.
\end{itemize}

% ---------------------------------------------------------
\subsection{AGN Accretion and Luminosity}

AGN feedback is governed by the accretion rate onto the central black hole:
\begin{equation}
\dot{M}_{\mathrm{BH}}
=
\epsilon_{\mathrm{BH}}
\frac{M_{\mathrm{gas}}}{t_{\mathrm{dyn}}},
\end{equation}
where $\epsilon_{\mathrm{BH}}$ is the accretion efficiency.

In RDCC, both $M_{\mathrm{gas}}$ and $t_{\mathrm{dyn}}$ are modified:
\begin{itemize}
\item reduced inflow lowers $M_{\mathrm{gas}}$,
\item delayed collapse increases $t_{\mathrm{dyn}}$,
\item reduced concentrations lower central densities.
\end{itemize}

The resulting accretion rate becomes
\begin{equation}
\dot{M}_{\mathrm{BH,RDCC}}
=
\dot{M}_{\mathrm{BH},\Lambda\mathrm{CDM}}
\left[
1
-
\chi_{\mathrm{BH}}
\left(
\frac{M_{\mathrm{rel}}}{M}
\right)^{\beta_{\mathrm{BH}}}
\right],
\end{equation}
with $\chi_{\mathrm{BH}} \sim 0.05$.

Consequences:
\begin{itemize}
\item AGN luminosities in low-mass halos are reduced,
\item AGN duty cycles become more intermittent,
\item the coupling efficiency of AGN energy increases.
\end{itemize}

% ---------------------------------------------------------
\subsection{AGN Energy Coupling Efficiency}

The AGN energy injection rate is
\begin{equation}
\dot{E}_{\mathrm{AGN}}
=
\epsilon_{f}
\dot{M}_{\mathrm{BH}} c^{2},
\end{equation}
where $\epsilon_{f}$ is the coupling efficiency.

In RDCC, shallower potentials and reduced gas densities increase the effective coupling:
\begin{equation}
\epsilon_{f,\mathrm{RDCC}}
=
\epsilon_{f,\Lambda\mathrm{CDM}}
\left[
1
+
\chi_{f}
\left(
\frac{M_{\mathrm{rel}}}{M}
\right)^{\beta_{f}}
\right],
\end{equation}
with $\chi_{f} \sim 0.05$.

Implications:
\begin{itemize}
\item AGN-driven winds propagate further,
\item heating is more spatially extended,
\item cold inflow streams are more easily disrupted.
\end{itemize}

% ---------------------------------------------------------
\subsection{Combined Feedback and Baryon Redistribution}

The combined effect of SN and AGN feedback in RDCC leads to:
\begin{itemize}
\item reduced central gas densities,
\item expanded CGM structures,
\item lower baryon fractions in halos with $M \lesssim 10^{12} M_{\odot}$,
\item smoother gas profiles,
\item enhanced entropy in the intracluster medium (ICM).
\end{itemize}

The baryon fraction becomes
\begin{equation}
f_{b,\mathrm{RDCC}}(M)
=
f_{b,\Lambda\mathrm{CDM}}(M)
\left[
1
-
\chi_{b}
\left(
\frac{M_{\mathrm{rel}}}{M}
\right)^{\beta_{b}}
\right],
\end{equation}
with $\chi_{b} \sim 0.1$.

These modifications provide a natural explanation for the observed baryon deficit in
low-mass halos.

% ---------------------------------------------------------
% SECTION 4 — THE RDCC BARYONIC CORRECTION MODEL FOR THE MATTER POWER SPECTRUM
% ---------------------------------------------------------
\section{The RDCC Baryonic Correction Model for the Matter Power Spectrum}

Baryonic processes such as gas cooling, star formation, and feedback redistribute matter
within halos and modify the matter power spectrum on small and intermediate scales. In the
standard $\Lambda$CDM framework, these effects are typically modeled using baryonic
correction models (BCMs) calibrated to hydrodynamical simulations. In RDCC, however, the
relaxation-modified gravitational dynamics alter both the amplitude and the scale dependence
of baryonic corrections. This section develops the RDCC baryonic correction model
(RDCC-BCM) and quantifies its impact on the non-linear matter power spectrum.

The RDCC-BCM incorporates:
\begin{itemize}
\item relaxation-modified halo profiles (Companion~XIX),
\item enhanced gas ejection due to reduced escape velocities,
\item weakened stellar contraction due to reduced star formation efficiency,
\item a smooth, universal suppression tied to the relaxation scale $k_{\mathrm{rel}}$.
\end{itemize}

% ---------------------------------------------------------
\subsection{Structure of the RDCC-BCM}

The RDCC-BCM modifies the matter power spectrum through three physically distinct
components:
\begin{enumerate}
\item dark matter redistribution due to relaxation-modified halo profiles,
\item gas ejection and expansion driven by enhanced SN and AGN feedback,
\item stellar contraction weakened by reduced star formation efficiency.
\end{enumerate}

The total correction factor is written as
\begin{equation}
B_{\mathrm{RDCC}}(k)
=
B_{\mathrm{grav}}(k)\,
B_{\mathrm{gas}}(k)\,
B_{\mathrm{stars}}(k).
\end{equation}

The full non-linear matter power spectrum becomes
\begin{equation}
P_{\mathrm{NL,RDCC+bar}}(k)
=
P_{\mathrm{NL,RDCC}}(k)\,
B_{\mathrm{RDCC}}(k),
\end{equation}
where $P_{\mathrm{NL,RDCC}}(k)$ is the relaxation-modified non-linear power spectrum
derived in Companion~XIX.

% ---------------------------------------------------------
\subsection{Dark Matter Redistribution Term}

Relaxation-modified halo profiles lead to a smooth suppression of small-scale power. This
effect is captured by
\begin{equation}
B_{\mathrm{grav}}(k)
=
1
-
A_{\mathrm{grav}}
\exp\!\left[
-
B_{\mathrm{grav}}
\left(
\frac{k}{k_{\mathrm{rel}}}
\right)^{2}
\right],
\end{equation}
with typical values
\begin{equation}
A_{\mathrm{grav}} \sim 0.05,
\qquad
B_{\mathrm{grav}} \sim 0.02.
\end{equation}

This term reproduces the RDCC halo-model suppression derived in Companion~XIX.

% ---------------------------------------------------------
\subsection{Gas Ejection Term}

Gas ejection due to SN and AGN feedback is enhanced in RDCC because of shallower
potentials and reduced inflow rates. The gas ejection radius scales as
\begin{equation}
r_{\mathrm{ej,RDCC}}
=
r_{\mathrm{ej},\Lambda\mathrm{CDM}}
\left[
1
+
\chi_{\mathrm{ej}}
\left(
\frac{M_{\mathrm{rel}}}{M}
\right)^{\beta_{\mathrm{ej}}}
\right],
\end{equation}
with $\chi_{\mathrm{ej}} \sim 0.1$.

The corresponding suppression factor is
\begin{equation}
B_{\mathrm{gas}}(k)
=
\exp\!\left[
-
A_{\mathrm{gas}}
\left(
\frac{k}{k_{\mathrm{ej}}}
\right)^{2}
\right],
\end{equation}
where $k_{\mathrm{ej}} \sim r_{\mathrm{ej}}^{-1}$.

This term dominates the baryonic suppression at
\begin{equation}
k \sim 1\text{--}10\,h\,\mathrm{Mpc}^{-1}.
\end{equation}

% ---------------------------------------------------------
\subsection{Stellar Contraction Term}

Stars contract the central potential and enhance small-scale power. In RDCC, reduced star
formation efficiency weakens this effect. The stellar contraction term is modeled as
\begin{equation}
B_{\mathrm{stars}}(k)
=
1
+
A_{\mathrm{stars}}
\left(
\frac{k}{k_{\mathrm{stars}}}
\right)^{2}
\exp\!\left[
-
B_{\mathrm{stars}}
\left(
\frac{k}{k_{\mathrm{stars}}}
\right)^{2}
\right],
\end{equation}
with
\begin{equation}
A_{\mathrm{stars,RDCC}}
=
A_{\mathrm{stars},\Lambda\mathrm{CDM}}
\left[
1
-
\chi_{\mathrm{stars}}
\left(
\frac{M_{\mathrm{rel}}}{M}
\right)^{\beta_{\mathrm{stars}}}
\right],
\end{equation}
where $\chi_{\mathrm{stars}} \sim 0.1$.

Consequences:
\begin{itemize}
\item stellar contraction is weaker than in $\Lambda$CDM,
\item small-scale enhancement is reduced,
\item central density profiles are smoother.
\end{itemize}

% ---------------------------------------------------------
\subsection{Combined RDCC Baryonic Suppression}

Combining the three components yields a smooth, physically motivated suppression of the
matter power spectrum. Typical values are:
\begin{itemize}
\item $\sim 5\%$ suppression at $k \sim 0.5\,h\,\mathrm{Mpc}^{-1}$,
\item $\sim 10$–$15\%$ suppression at $k \sim 1\,h\,\mathrm{Mpc}^{-1}$,
\item $\sim 20$–$30\%$ suppression at $k \sim 5\,h\,\mathrm{Mpc}^{-1}$.
\end{itemize}

These values are consistent with hydrodynamical simulations once relaxation effects are
included.

% ---------------------------------------------------------
\subsection{Comparison with Other Small-Scale Models}

The RDCC-BCM differs fundamentally from other small-scale suppression models:
\begin{itemize}
\item warm dark matter introduces a sharp cutoff, absent in RDCC,
\item self-interacting dark matter modifies halo cores but not the power spectrum smoothly,
\item early dark energy enhances small-scale power rather than suppressing it.
\end{itemize}

The RDCC suppression is smooth, universal, and tied directly to the relaxation scale
$k_{\mathrm{rel}}$.

% ---------------------------------------------------------
% SECTION 5 — BARYON REDISTRIBUTION IN HALOS AND THE INTERGALACTIC MEDIUM
% ---------------------------------------------------------
\section{Baryon Redistribution in Halos and the Intergalactic Medium}

The distribution of baryons within halos and the intergalactic medium (IGM) is shaped by
the interplay between gravitational collapse, gas cooling, star formation, and feedback. In
the standard $\Lambda$CDM framework, baryons are redistributed primarily through supernova
(SN) and active galactic nuclei (AGN) feedback. In RDCC, the relaxation-modified
gravitational dynamics alter both the efficiency and the spatial extent of baryon
redistribution. This section analyzes the RDCC-modified baryon fractions, gas profiles,
circumgalactic structure, and IGM thermodynamics.

Relaxation affects baryon redistribution through:
\begin{itemize}
\item reduced escape velocities in halos with $M < M_{\mathrm{rel}}$,
\item enhanced SN- and AGN-driven outflows,
\item reduced inflow and smoother central gas densities,
\item more extended CGM structures,
\item modified IGM heating and enrichment.
\end{itemize}

% ---------------------------------------------------------
\subsection{Baryon Fractions in RDCC Halos}

The baryon fraction of a halo is defined as
\begin{equation}
f_{b}(M)
=
\frac{M_{\mathrm{baryon}}}{M_{\mathrm{halo}}}.
\end{equation}

In RDCC, enhanced outflows and reduced inflow lead to
\begin{equation}
f_{b,\mathrm{RDCC}}(M)
=
f_{b,\Lambda\mathrm{CDM}}(M)
\left[
1
-
\chi_{b}
\left(
\frac{M_{\mathrm{rel}}}{M}
\right)^{\beta_{b}}
\right],
\end{equation}
with $\chi_{b} \sim 0.1$ and $\beta_{b} \sim 2/3$.

Consequences:
\begin{itemize}
\item dwarf galaxies lose a larger fraction of their baryons,
\item galaxy groups show reduced hot-gas fractions,
\item clusters remain largely unaffected ($M \gg M_{\mathrm{rel}}$).
\end{itemize}

This naturally explains the observed baryon deficit in low-mass halos.

% ---------------------------------------------------------
\subsection{Gas Density and Temperature Profiles}

The gas density profile is typically modeled as a $\beta$-profile:
\begin{equation}
\rho_{\mathrm{gas}}(r)
=
\rho_{0}
\left[
1
+
\left(
\frac{r}{r_{c}}
\right)^{2}
\right]^{-3\beta_{\mathrm{gas}}/2}.
\end{equation}

In RDCC, reduced concentrations increase the core radius $r_{c}$, and enhanced feedback
lowers the central density $\rho_{0}$. We model this as
\begin{equation}
r_{c,\mathrm{RDCC}}
=
r_{c,\Lambda\mathrm{CDM}}
\left[
1
+
\chi_{c}
\left(
\frac{M_{\mathrm{rel}}}{M}
\right)^{\beta_{c}}
\right],
\end{equation}
with $\chi_{c} \sim 0.05$.

The temperature profile is similarly modified:
\begin{equation}
T_{\mathrm{RDCC}}(r)
=
T_{\Lambda\mathrm{CDM}}(r)
\left[
1
-
\chi_{T}
\left(
\frac{M_{\mathrm{rel}}}{M}
\right)^{\beta_{T}}
\right],
\end{equation}
with $\chi_{T} \sim 0.05$.

Implications:
\begin{itemize}
\item gas in low-mass halos is cooler,
\item temperature gradients are smoother,
\item central X-ray luminosities are reduced.
\end{itemize}

% ---------------------------------------------------------
\subsection{Circumgalactic Medium (CGM) Structure}

The CGM is highly sensitive to feedback and gravitational potential depth. In RDCC:
\begin{itemize}
\item SN-driven winds propagate further,
\item AGN-driven bubbles expand more efficiently,
\item cold inflow streams are disrupted by vorticity,
\item the CGM becomes more extended and less clumpy.
\end{itemize}

The CGM density profile becomes
\begin{equation}
\rho_{\mathrm{CGM,RDCC}}(r)
=
\rho_{\mathrm{CGM},\Lambda\mathrm{CDM}}(r)
\left[
1
-
\chi_{\mathrm{CGM}}
\left(
\frac{M_{\mathrm{rel}}}{M}
\right)^{\beta_{\mathrm{CGM}}}
\right],
\end{equation}
with $\chi_{\mathrm{CGM}} \sim 0.1$.

Consequences:
\begin{itemize}
\item lower HI column densities,
\item reduced Lyman-$\alpha$ absorption,
\item more extended OVI/OVIII halos.
\end{itemize}

% ---------------------------------------------------------
\subsection{Gas Ejection and Redistribution Radius}

The radius to which gas is ejected by feedback is
\begin{equation}
r_{\mathrm{ej}}
\sim
\frac{E_{\mathrm{fb}}}{|\Phi_{\mathrm{halo}}|},
\end{equation}
where $E_{\mathrm{fb}}$ is the feedback energy.

In RDCC:
\begin{itemize}
\item $E_{\mathrm{fb}}$ increases (enhanced SN and AGN efficiency),
\item $|\Phi_{\mathrm{halo}}|$ decreases (delayed collapse).
\end{itemize}

Thus,
\begin{equation}
r_{\mathrm{ej,RDCC}}
=
r_{\mathrm{ej},\Lambda\mathrm{CDM}}
\left[
1
+
\chi_{\mathrm{ej}}
\left(
\frac{M_{\mathrm{rel}}}{M}
\right)^{\beta_{\mathrm{ej}}}
\right],
\end{equation}
with $\chi_{\mathrm{ej}} \sim 0.2$.

Consequences:
\begin{itemize}
\item stronger baryon evacuation in dwarfs,
\item extended warm-hot gas in groups,
\item mild effects in clusters.
\end{itemize}

% ---------------------------------------------------------
\subsection{IGM Thermodynamics}

The IGM temperature–density relation is
\begin{equation}
T
=
T_{0}
\Delta^{\gamma - 1},
\end{equation}
where $\Delta = \rho/\bar{\rho}$.

In RDCC:
\begin{itemize}
\item reduced small-scale power lowers shock heating,
\item enhanced vorticity increases turbulent heating,
\item gas ejection enriches the IGM with metals.
\end{itemize}

We model the modified normalization as
\begin{equation}
T_{0,\mathrm{RDCC}}
=
T_{0,\Lambda\mathrm{CDM}}
\left[
1
-
\chi_{1}
+
\chi_{2}
\right],
\end{equation}
with $\chi_{1} \sim 0.05$ (reduced shocks) and $\chi_{2} \sim 0.02$ (turbulence).

Implications:
\begin{itemize}
\item the low-density IGM is slightly cooler,
\item the temperature distribution is broader,
\item the warm-hot intergalactic medium (WHIM) fraction increases.
\end{itemize}

% ---------------------------------------------------------
\subsection{Observational Signatures}

RDCC predicts:
\begin{itemize}
\item reduced X-ray luminosities in low-mass halos,
\item enhanced SZ pressure profiles in groups,
\item lower HI column densities in the CGM,
\item more extended OVI/OVIII halos,
\item increased WHIM fraction,
\item smoother gas density profiles,
\item suppressed baryon fractions in dwarfs and groups.
\end{itemize}

These signatures are testable with:
\begin{itemize}
\item X-ray: Athena, XRISM,
\item SZ: SPT, ACT, Simons Observatory,
\item HI: SKA, MeerKAT,
\item UV absorption: HST/COS, LUVOIR,
\item Lyman-$\alpha$: DESI, WEAVE.
\end{itemize}

% ---------------------------------------------------------
% SECTION 6 — CONCLUSION
% ---------------------------------------------------------
\section{Conclusion}

This Companion has developed the first comprehensive framework for baryonic physics in the
Relaxation-Driven Cyclic Cosmology (RDCC). Building on the scale-dependent growth derived
in Companion~XVIII and the non-linear structure formation analysis of Companion~XIX, we
have shown that relaxation-induced modifications of gravitational dynamics directly influence
gas cooling, star formation, feedback energetics, and the redistribution of baryons within halos
and the intergalactic medium.

Delayed halo collapse and reduced gravitational potentials lower virial temperatures,
increase cooling times, and suppress early star formation. Relaxation-induced vorticity modifies
filamentary inflow and reduces the gas supply available for star formation, while enhanced SN
and AGN feedback efficiency drives more extended outflows and redistributes baryons to larger
radii. These effects naturally produce lower baryon fractions in low-mass halos, smoother gas
profiles, and an enhanced warm-hot intergalactic medium (WHIM) fraction.

We constructed the RDCC baryonic correction model (RDCC-BCM), which combines
relaxation-modified halo profiles, enhanced gas ejection, and weakened stellar contraction.
The resulting suppression of the matter power spectrum is smooth, physically motivated, and
tied directly to the relaxation scale $k_{\mathrm{rel}}$. This distinguishes RDCC from
$\Lambda$CDM even in the presence of baryonic uncertainties and provides a robust target for
upcoming surveys.

Overall, the baryonic predictions developed in this Companion demonstrate that RDCC
offers a coherent and testable framework for understanding the interplay between gravitational
relaxation and baryonic physics. The model naturally suppresses small-scale power, modifies
halo gas fractions, and reshapes the circumgalactic and intergalactic media, all while remaining
consistent with current observations. Future missions such as Euclid, LSST, SKA, Athena, and
XRISM will be able to test these predictions with unprecedented precision.

This Companion completes the extension of RDCC into the baryonic sector and establishes
a direct observational bridge between relaxation dynamics and the astrophysics of galaxy
formation.


\end{document}
